% Generated draft: requires mathematical and source review.
\documentclass[11pt]{article}
\usepackage[margin=25mm]{geometry}
\usepackage{fontspec}
\setmainfont{Latin Modern Roman}
\usepackage{amsmath,amssymb,mathtools,unicode-math}
\setmathfont{Latin Modern Math}
\usepackage{enumitem,needspace,xurl,hyperref}
\setlength{\parindent}{0pt}
\setlength{\parskip}{5pt}
\setlength{\emergencystretch}{3em}
\newcommand{\N}{\mathbb{N}}
\newcommand{\Z}{\mathbb{Z}}
\newcommand{\Q}{\mathbb{Q}}
\newcommand{\R}{\mathbb{R}}
\newcommand{\C}{\mathbb{C}}
\newcommand{\F}{\mathbb{F}}
\newcommand{\A}{\mathbb{A}}
\newcommand{\PP}{\mathbb{P}}
\newcommand{\E}{\mathbb{E}}
\newcommand{\eps}{\varepsilon}
\newcommand{\dd}{\,\mathrm{d}}
\begin{document}
\section*{UnsolvedMath ID: 800836}
\textbf{Source problem code:} catalogue entry
\subsection{Definitions and mathematical statement}
The supplied statement refers to a mechanism and an apparatus described in preceding source material that is absent from this record. It also introduces parameters \(n_i,m_i,r_i\), but supplies neither their definitions nor their domains, index set, or relationships. The terms ``cascade'' and ``real hydrodynamics'' are not mathematically specified in the supplied record.

With these unresolved references retained, the source asks:
\begin{enumerate}
\item Does there exist a mechanism describing a cascade in real hydrodynamics that is analogous to the mechanism referred to in the preceding source discussion?
\item Is it possible to define the parameters \(n_i,m_i,r_i\) so that the apparatus referred to in that discussion works as intended?
\end{enumerate}
The governing equations, admissible hydrodynamic setting, relevant boundary conditions, meaning of a cascade, comparison mechanism, and required properties of the apparatus are not supplied. Consequently, these questions cannot yet be expressed as a fully specified mathematical existence statement with explicit domains, hypotheses, and quantifiers.
\subsection{Short English statement}
The problem asks whether a previously described mechanism has an analogue for hydrodynamic cascades, and whether parameters \(n_i,m_i,r_i\) can be defined to make the associated apparatus work. The missing preceding discussion prevents a precise mathematical formulation.
\subsection{Sources}
\begin{itemize}
\item Local UnsolvedMath registry: permanent ID \(800836\), source problem code \texttt{catalogue entry}, title \emph{Geometric Hydrodynamics in Open Problems --- Problem 35}. The supplied record attributes the problem to Khesin, Misiolek, and Modin, with proposed year 2022, and labels it open; this label is supplied registry metadata, not an independently verified mathematical status.
\item Khesin, Misiolek, and Modin, \emph{Geometric Hydrodynamics in Open Problems} (2022), Problem 35. Supplied source URL: \url{https://arxiv.org/abs/2205.01143}. The record reports access on 2026-10-16 and extraction by \texttt{arxiv\_latex}; these are supplied provenance details, not a claim of independent source consultation.
\item Supplied upstream revision: \texttt{d371ee917ac9647fbb06cd6fc529deec6f3dda82}.
\end{itemize}
\end{document}
